\chapter{Python 入门与完整逐行注释模板}
\section{先看懂程序的几个基本规则}
\texttt{import} 是加载工具箱；\texttt{as np} 是为 NumPy 取简称；\texttt{def} 定义可反复使用的函数；冒号后的缩进表示函数或循环的内部。\texttt{return} 把计算结果交给调用者。Python 区分大小写。

单个等号是赋值；两个等号是判断相等。\texttt{**2} 是平方；\texttt{*} 是乘法；\texttt{/} 是除法。数学中的并列相乘在代码里必须写乘号。\texttt{\#} 后的文字是给人看的解释，不参与运算。\texttt{assert} 后面的条件不成立时就停止程序，避免把失败结果当答案。

数组下标从零开始，\texttt{y[0]} 表示第一个分量，\texttt{y[-1]} 表示最后一个。切片 \texttt{u[1:-1]} 表示去掉两端；\texttt{u[:-2]} 去掉最后两个；\texttt{u[2:]} 去掉最前两个，它们可分别与内部点的左邻居、右邻居对应。

\texttt{np.array} 创建数值数组；\texttt{np.linspace(a,b,n)} 在含端点的区间生成 $n$ 个等距点，其中 $a,b$ 为端点、$n$ 为点数；\texttt{np.exp} 与 \texttt{np.log} 是指数与自然对数；\texttt{np.abs} 是绝对值；\texttt{np.max} 是最大值；\texttt{np.all} 要求全部元素都满足条件。\texttt{lambda z: ...} 是短函数写法，冒号后是返回的表达式。

\section{运行步骤}
安装 Python 3.11 或以上版本后，在本手册目录打开终端。先运行安装命令，再运行相应文件。命令不是在 Python 的三个大于号交互提示符后输入，而是在系统终端输入。下面每行的井号后都解释该行目的。
\begin{Verbatim}[breaklines=true,breakanywhere=true,fontsize=\small]
python -m pip install numpy scipy  # 安装三个数值示例所需的数组与科学计算库
python code/ode_template.py  # 求解冷却ODE并核验解析解与事件时间
python code/heat_1d.py  # 比较三种热差分方法并检查网格加密误差
python code/optimize_template.py  # 运行带约束优化和冷却参数拟合
\end{Verbatim}
网络不可用时，需先在有网络的环境准备匹配 Python 版本的安装包。模板自身不访问网络，也不需要 API 密钥。输出中的误差允许末尾数字随平台略变；是否通过断言比完全逐字相同更重要。

\section{ODE：冷却到指定温度}
输入为环境温度、初始温度、冷却系数与目标温度；输出为误差和到达目标的时间。替换方程时同时修改状态初值、参数单位和事件函数。\texttt{rtol} 是相对误差控制，\texttt{atol} 是绝对误差控制；它们控制内部误差估计，不是最终误差的严格保证。
\VerbatimInput[breaklines=true,breakanywhere=true,fontsize=\small,numbers=left,numbersep=5pt]{code/ode_template.py}

\section{一维热传导：显式、隐式与 Crank--Nicolson}
本模板只实现均匀材料、常热扩散率、无热源、两端恒为零的教学验证问题。温度也可理解为相对于固定端点温度的超温。它不是可直接承接任意对流边界、多层材料的通用成品；改变这些物理条件时须按正文重新装配系数与右端项。
\VerbatimInput[breaklines=true,breakanywhere=true,fontsize=\small,numbers=left,numbersep=5pt]{code/heat_1d.py}

\section{非线性优化与参数拟合}
第一个例子有手算最优解；第二个例子用已知参数生成无噪声数据，检查能否找回参数。实际有噪声数据不会精确恢复真值，还应增加置信区间、独立验证与可辨识性分析。
\VerbatimInput[breaklines=true,breakanywhere=true,fontsize=\small,numbers=left,numbersep=5pt]{code/optimize_template.py}

\chapter{论文落地与查阅路线}
\section{按问题找工具}
\begin{longtable}{p{0.24\textwidth}p{0.31\textwidth}p{0.35\textwidth}}
\toprule 题目现象 & 起步模型 & 必须核查 \\ \midrule
温度随时间变 & 集总热容 ODE & Biot 数、环境是否恒定 \\
物体内部温差明显 & 热传导 PDE 与差分 & 边界、热源、材料界面、收敛 \\
振动、轨迹、运动 & 牛顿方程与一阶系统 & 力的方向、坐标、接触事件 \\
管道输水、压差 & 连续方程与能量方程 & 雷诺数、损失、适用流态 \\
光强、成像、干涉 & 对应几何光学或波动光学 & 波长、角度、相干性与单位 \\
已知曲线想求参数 & 非线性最小二乘 & 误差结构、参数相关性 \\
选尺寸、功率、时间 & 带约束非线性优化 & 可行性、多初值、物理限制 \\
重复试验或测量噪声 & 概率模型与统计推断 & 独立性、分布假设、样本量 \\
数据多但机理不全 & SINDy、PINN 等选读 & 先有基线、数据覆盖与验证成本 \\
\bottomrule
\end{longtable}

\section{把公式写进论文的模板}
可以这样写：“将研究对象视为……，忽略……，依据……守恒建立……方程。式中各符号与单位列于符号表；初始条件由……获得，边界条件对应……。参数采用……数据估计。使用……方法求解，并通过解析特例、网格加密与独立观测检验。最终输出……，适用于……范围。”省略号必须由题目事实补足。

不要把代码运行成功写成“模型正确”；不要把一次优化终止写成“全局最优”；不要把文献中的某个材料参数直接搬到另一材料；不要把未经验证的前沿模型包装成必然优于经典方法。

\section{交付文件说明}
主文档是 PDF；\texttt{main.tex} 是可编辑 LaTeX 源文件；\texttt{chapters} 内保存章节；\texttt{code} 内是可独立运行的完整模板。公式均保留为 LaTeX 数学表达式，可以复制到论文。所有公式的符号含义以紧邻正文为准；无量纲比例和编号也已说明。

\section{基础文献的进一步阅读}
\begin{enumerate}
\item Boyd, S.; Vandenberghe, L. \emph{Convex Optimization}. Cambridge University Press, 2004. 凸性、对偶与 KKT。作者开放教材：\url{https://web.stanford.edu/~boyd/cvxbook/}。
\item NIST/SEMATECH. \emph{e-Handbook of Statistical Methods}. 统计分布、拟合、诊断与区间。\url{https://www.itl.nist.gov/div898/handbook/}。
\item SciPy 官方文档。实现时以实际安装版本为准，重点查 \texttt{solve\_ivp}、\texttt{solve\_banded}、\texttt{minimize}、\texttt{least\_squares}。\url{https://docs.scipy.org/doc/scipy/reference/}。
\item OpenStax. \emph{University Physics}, Volumes 1--3. 基础力学、流体、热学、光学。\url{https://openstax.org/subjects/science}。各物理章节另列具体入口。
\end{enumerate}
本文为教学性综合整理，文献观点用自己的语言概括；方法公式不是其实际性能的保证。网络来源核查日期为 2026 年 9 月 10 日；前沿章节是代表性路线，未声称穷尽或代表最新全部进展。
